\chapter{Related Work and Positioning}\label{ch:m-related}

\section{GPU Communication Characterization}

Surveys of the xCCL ecosystem compare supported collectives, algorithms, and vendor scope, while topology studies show that nominally identical GPU operations can differ substantially with link placement and routing \cite{weingram_xccl_2023,li_evaluating_2020}. PICO provides controlled algorithm selection, profiling, and environment recording across MPI and NCCL. Its MPI allreduce sweeps include configurations where the default has approximately five times the latency of the best exposed alternative. PICO also goes beyond microbenchmarks: ATLAHS replay of LLM training traces projects runtime reductions of up to 44\% from benchmark-informed collective choices. This is simulation-based application assessment, not a measured training-run speedup \cite{pasqualoni_pico_2026}.

TuCCL combines algorithm and runtime-parameter optimization and reports 1.22--2.52-fold end-to-end training speedups for its evaluated workloads and comparison systems \cite{li_tuccl_2025}. Choi et al. likewise connect communication simulation and GPU kernel models to end-to-end application behavior \cite{choi_end--end_2020}. These studies establish that configuration and application context belong in the analysis. This thesis builds on that premise with a different evidence chain: a common CUDA/SYCL comparison including SHMEM paths, concrete oneCCL extensions, and measured execution of complete sparse CG solvers. It does not claim to introduce benchmark-to-application validation itself.

\section{GPU-Initiated and One-Sided Communication}

Agostini et al. study GPUDirect Async through microbenchmarks and a performance model, examining potential benefits in HPGMG-FV and CoMD-CUDA and reporting a BFS case without an advantage \cite{agostini_gpudirect_2018}. This is an important precedent for the thesis's conditional question. The mechanism, workloads, and measurement boundaries differ from later NVSHMEM studies, so a no-benefit BFS case there does not contradict benefits reported for other GPU-initiated implementations. NVSHMEM analyses emphasize symmetric memory, granularity, operation scope, and host-on-stream versus device-side behavior \cite{unat_landscape_2026,ma_demystifying_2026}.

Hamidouche et al.'s CCGrid study describes kernel-initiated networking through NCCL GIN and its DeepEP integration, including GDAKI and proxy backends \cite{hamidouche_kernel-initiated_2026}. The earlier NCCL GIN preprint provides closely related architectural context \cite{hamidouche_gpu-initiated_2025}; these are not treated as independent replications. Shen et al. study low-latency GPU collectives using symmetric-memory and synchronization techniques, with hardware features that are not all present on Leonardo \cite{shen_every_2026}. The new intra-node NCCL summary in this thesis motivates a follow-up on A100 but does not reproduce those papers' device-kernel or direct-network experiments.

This thesis complements these studies by comparing persistent device-initiated NVSHMEM neighbor kernels with host-initiated stacks under one pattern matrix, using OSHMPI as a host-driven OpenSHMEM path, and implementing SHMEM-family support beneath oneCCL at two readiness levels. The complete OSHMPI backend supplies benchmark evidence; the NVSHMEM foundation supplies runtime and staging validation. Neither is a claim to have introduced one-sided or GPU-initiated communication.

\section{Communication in Conjugate Gradient}

Trotter et al. develop the aCG suite and compare host- and device-initiated MPI, NCCL/RCCL, and NVSHMEM strategies on large NVIDIA and AMD systems \cite{trotter_cpu-_2025}. Their results show that a monolithic device-side solver can remove CPU synchronization but does not dominate every host-initiated variant at every scale. This thesis uses the same application family and reporting convention, adds an OSHMPI communicator and a SYCL implementation, and asks a different methodological question: whether a prior pattern-oriented benchmark would have predicted the application result.

\section{Position of This Work}

The contribution is the specific combination of a pattern-oriented CUDA/SYCL comparison, baseline-relative oneCCL extensions, and measured validation through OSHMPI and SYCL extensions to aCG on Leonardo. Existing work already connects mechanisms, configuration, and application performance through modeling, simulation, and execution. Here the shared evidence chain makes it possible to identify which primitives transfer quantitatively and why an application ranking changes. The five contributions in Table~\ref{tab:m-artifact-evidence} provide the software and measurements for that analysis, while the preliminary NCCL comparison identifies a further candidate for validation. Claims of novelty concern this implementation and evaluation scope, rather than a first use of GPU initiation or a first connection between benchmarks and applications.

